\documentclass[a4paper]{article}

%%% To be incput by template-*.tex

\usepackage[left=20mm,right=20mm,top=20mm,bottom=20mm]{geometry}
\usepackage{amsmath, amssymb, amsthm}
\usepackage{tikz}
\usetikzlibrary{arrows}
\tikzset{every node/.style={draw, circle, inner sep=3pt}}
\usepackage{enumitem}
\usepackage{systeme}
\usepackage{multicol}
\usepackage{siunitx}
\usepackage{qrcode}
\usepackage{bm}

\pagenumbering{gobble}

\usepackage{CJKutf8}
\newcommand{\CJKEN}[1]{\begin{CJK}{UTF8}{bkai}#1\end{CJK}}

\newcommand{\rboxed}[1]{{\colorlet{oldcolor}{.}\color{red}\boxed{\color{oldcolor}#1}}}

%\parindent=0pt



\begin{document}

%%%FOOTER: KEY, CHECKCODE, MIDDLE TEXT
\begin{tikzpicture}[remember picture, overlay]

\node[draw=none, rectangle] at ([shift={(2cm,2cm)}]current page.south west){
\qrcode[height=2cm,level=H
%version=10 % 1 ~ 40, lines = 16 + 4 v
]{
%%%%%KEY
}
};
\node[draw=none, rectangle, above] at ([shift={(2cm,3cm)}]current page.south west) {
%%%%%KEY
};

\draw[very thick] ([shift={(-3cm,1cm)}]current page.south east) rectangle ([shift={(-1cm,3cm)}]current page.south east);
\node[draw=none, rectangle, above] at ([shift={(-2cm,3cm)}]current page.south east) {check code};
\node[draw=none, rectangle, scale=2, transform shape] at ([shift={(-2cm,2cm)}]current page.south east){
\Huge
%%%%%CHECKCODE
};

\end{tikzpicture} 

\begin{Large}
\centerline{\CJKEN{姓名} Name : \underline{\hspace{10em}} \hfill \CJKEN{學號} Student ID \# : \underline{\hspace{10em}}}
%%%%%HEADING
%%% Sample heading
% \centerline{\CJKEN{國立中山大學} \hfill NATIONAL SUN YAT-SEN UNIVERSITY}
%\vspace{0.5cm}
%\centerline{\textbf{\CJKEN{線性代數（一）}} \hfill \textbf{MATH 103 / GEAI 1215: Linear Algebra I}}
%\vspace{0.5cm}
%\centerline{\textbf{\CJKEN{第一次期中考}} \hfill October 8, 2018 \hfill \textbf{Midterm 1}}
%\vspace{0.5cm}

\vspace{1cm}
%%%%%QUESTION

\vspace{1cm}
\noindent\textbf{Solution.}

%%%%%SOLUTION


\vfill
\begin{center}
\begin{obeylines}
Indicating your answer by \textbf{underlining it} or \textbf{circling it}.

Compute the \textbf{check code} and fill it into the \textbf{box on the right}.
\end{obeylines}
\end{center}


\end{Large}
\end{document}
